\chapter{Konzeptionelle Grundlagen}
\section{Technische Grundlagen}
\subsection{Containerarchitektur}
Containerisierung bezeichnet die Kapselung von Anwendungen samt ihrer Laufzeitabhängigkeiten in isolierte, portable Einheiten – sogenannte Container. Im Unterschied zu klassischen virtuellen Maschinen teilen sich Container den Kernel des Host-Betriebssystems, was einen erheblich geringeren Ressourcenoverhead sowie kürzere Startzeiten bedingt (vgl. Merkel 2014, S. 2 f.). Docker hat sich als De-facto-Standard für die Erstellung und den Betrieb von Containern etabliert, indem Laufzeitumgebungen durch deklarative Dockerfiles vollständig reproduzierbar definiert werden (vgl. Anderson 2015, S. 1 f.).

Für Multi-Container-Anwendungen erweitert Docker Compose diesen Ansatz um eine zentrale, deklarative Konfigurationsdatei zur Vernetzung und Steuerung mehrerer Dienste (vgl. Stoneman 2020, S. 14 ff.). In Verbindung mit dem Microservice-Paradigma – bei dem funktionale Einheiten in lose gekoppelte, unabhängig deploybare Dienste zerlegt werden – entfaltet Containerisierung ihr volles Potenzial: Einzelne Komponenten lassen sich bedarfsgerecht skalieren, unabhängig aktualisieren und replizieren, ohne das Gesamtsystem zu beeinträchtigen (vgl. Newman 2015, S. 1 ff.).

\subsection{Anwendung von Python}
Python bildet aufgrund seiner lesbaren Syntax und seines umfangreichen Bibliotheksökosystems die technologische Basis der vorliegenden Systemarchitektur (vgl. VanderPlas 2016, S. 2 f.). Die Datenverarbeitung stützt sich auf pandas, dessen DataFrame-Konzept tabellarische Daten effizient vorhält und für Einlesen, Bereinigung sowie Aggregation der Rohdaten eingesetzt wird (vgl. McKinney 2017, S. 97 ff.). Das spaltenorientierte Apache-Parquet-Format wird über PyArrow eingebunden und bietet signifikante Lese- und Speichervorteile gegenüber zeilenorientierten Formaten wie CSV (vgl. Vohra 2016, S. 11 f.). Für die Imputation fehlender Werte kommt der IterativeImputer aus scikit-learn zum Einsatz, der das MICE-Verfahren implementiert (vgl. Pedregosa et al. 2011, S. 2825 f.).

Die statistische Modellierung stützt sich auf statsmodels, das Generalisierte Lineare Modelle (GLM) und Kleinstquadrateschätzung (OLS) über eine patsy-basierte Formelschnittstelle bereitstellt (vgl. Seabold/Perktold 2010, S. 93). Die HTTP-Schnittstellen zwischen den Containerdiensten werden über Flask in Kombination mit dem WSGI-Server Waitress realisiert (vgl. Grinberg 2018, S. 4 ff.). Die Visualisierungsschicht basiert auf Streamlit und Plotly Express, die interaktive Webanwendungen in reinem Python ermöglichen (vgl. Tamminga 2022, S. 1 f.).

\subsection{REST-API}
Representational State Transfer (REST) bezeichnet einen von Roy Thomas Fielding formalisierten Architekturstil für netzwerkbasierte Softwareschnittstellen, der auf dem HTTP-Protokoll aufsetzt (vgl. Fielding 2000, S. 76 ff.). Kernprinzipien sind strikte Client-Server-Trennung, Zustandslosigkeit sowie ein einheitliches Interface: Ressourcen werden über stabile URIs adressiert und über standardisierte HTTP-Methoden manipuliert, wobei Antworten typischerweise im JSON-Format übermittelt werden (vgl. Richardson/Ruby 2007, S. 7 ff.). Im Kontext von Microservice-Architekturen fungieren REST-APIs als primärer Kommunikationskanal zwischen Diensten und gewährleisten so funktionale Isolation sowie horizontale Skalierbarkeit (vgl. Newman 2015, S. 1 ff.).

\section{Vorstellung des Datensatzes}
Die empirische Grundlage dieser Arbeit bildet das Open-Data-Repertoire der New York City Taxi \& Limousine Commission (TLC). Der Datensatz umfasst detaillierte Aufzeichnungen über Millionen von Individualverkehrsbewegungen im Stadtgebiet New Yorks und gilt aufgrund seiner hohen zeitlichen sowie räumlichen Auflösung als Standardreferenz in der raumzeitlichen Datenanalyse. Herangezogen wird primär das Segment der „Yellow Taxi Data", welches den regulären Taxibetrieb im Stadtzentrum abbildet.

Strukturell handelt es sich um hochgradig granulare Transaktionsdaten: Jeder Datensatz repräsentiert eine einzelne Beförderung und beinhaltet Zeitstempel für Aufnahme und Absetzen der Fahrgäste, geographische Taxizonen sowie die zurückgelegte Fahrtstrecke. Ergänzend sind Tarifstrukturen (Fahrpreis, Steuern, Zuschläge) und die gewählte Zahlungsart erfasst. Angesichts eines Volumens von mehreren Gigabyte pro Jahr erfüllt der Datensatz die Charakteristika von Big Data. Die Vorverarbeitung erfordert dabei eine systematische Behandlung von Ausreißern – etwa unplausible Fahrtdauern oder fehlerhafte Tarifangaben – sowie vereinzelter Inkonsistenzen. Durch die Kombination aus zeitlicher Dichte und geographischer Präzision lassen sich komplexe urbane Mobilitätsmuster extrahieren, Nachfragespitzen identifizieren und ökonomische Trends ableiten.

\section{Angewendete Methoden}
\subsection{Multiple Imputation}
Fehlende Werte stellen in empirischen Datensätzen ein systematisches Problem dar, das bei naiver Behandlung zu verzerrten Schätzungen und einem Verlust statistischer Effizienz führt. Der listenweise Ausschluss unvollständiger Beobachtungen setzt voraus, dass Daten vollständig zufällig fehlen (MCAR, Missing Completely at Random) – eine Annahme, die in der Praxis selten zutrifft (vgl. Little/Rubin 2002, S. 41 ff.).

Als methodisch überlegene Alternative erzeugt die Multiple Imputation mehrere plausible Datensätze, die jeweils unterschiedliche Imputationswerte enthalten und damit die Unsicherheit über den wahren, nicht beobachteten Wert explizit abbilden (vgl. Rubin 1987, S. 75 ff.). Die leistungsfähige Ausprägung MICE (Multiple Imputation by Chained Equations) spezifiziert für jede unvollständige Variable ein separates Regressionsmodell und durchläuft diese Kette konditionaler Modelle iterativ bis zur Konvergenz (vgl. van Buuren/Groothuis-Oudshoorn 2011, S. 7 ff.). In scikit-learn ist dieses Prinzip als IterativeImputer implementiert (vgl. Pedregosa et al. 2011, S. 2826).

\subsection{Generalisiertes Lineares Modell}
Die klassische lineare Regression setzt eine kontinuierliche, normalverteilte Zielvariable voraus. Für Zähldaten oder strikt positive Mengengrößen sind diese Annahmen verletzt, was zu ineffizienten Schätzern führt (vgl. McCullagh/Nelder 1989, S. 1 ff.). Das von Nelder und Wedderburn 1972 eingeführte Generalisierte Lineare Modell (GLM) überwindet diese Einschränkungen durch drei Komponenten: einen linearen Prädiktor, eine Verknüpfungsfunktion sowie eine Verteilungsannahme aus der Exponentialfamilie (vgl. Nelder/Wedderburn 1972, S. 370 ff.). Für Zähldaten eignet sich die Poisson-Verteilung; bei Überdispersion bietet die Negativ-Binomialverteilung eine robustere Alternative (vgl. Cameron/Trivedi 2013, S. 75 f.). Für strikt positive, stetige Zielgrößen stehen Gamma- sowie Invers-Gaußsche Verteilung zur Verfügung (vgl. McCullagh/Nelder 1989, S. 292 ff.). Die Parameterschätzung erfolgt über die Maximum-Likelihood-Methode; zur Modellbewertung dienen das Akaike-Informationskriterium (AIC) sowie die p-Werte der Koeffizientenschätzungen (vgl. Fahrmeir/Tutz 2001, S. 86 ff.).
